% [原文 p.179]
对任意 $X\subset\mathbb{A}^n$，$B_0^wX$ 定义为 $X$ 在 $B_0^w\mathbb{A}^n$ 中的双有理变换。\hfill$\square$
% [Tutor Note: 上段是定义 4.56 的末段，起于 p.179 而定义本体在 p.178；本块按页界承接，故以普通段落收尾，不重复“定义 4.56”的标题与编号。]

下面证明本节的主要结果。

\begin{theorem*}{4.57}
（[Lau77, Rei76]）设 $(0\in X)$ 是二维椭圆奇点，$f:Y\to X$ 是极小奇点消解。有效例外除子 $Z$ 由 $\omega_Y=f^*\omega_X(-Z)$ 确定。记 $k=-(Z\cdot Z)$。
\begin{enumerate}[label=(\arabic*),leftmargin=2.4em,itemsep=1pt,topsep=2pt]
\item 假设 $k\ge 3$。此时 $(0\in X)$ 的重数为 $k$、嵌入维数为 $k$。任取嵌入 $(0\in X)\hookrightarrow(0\in\mathbb{A}^k)$，$x_i$ 是 $\mathbb{A}^k$ 的坐标系，$w$ 是由 $w(x_1,\dots,x_k)=(1,\dots,1)$ 确定的权。
\item 设 $k=2$。此时 $(0\in X)$ 的重数为 2、嵌入维数为 3。适当重取解析坐标后，$X$ 由如下形式的方程定义：
\[ z^2+q(x,y)=0\qquad(\mult_0q\ge4). \]
$w$ 是由 $w(x,y,z)=(1,1,2)$ 确定的权。
\item 设 $k=1$。此时 $(0\in X)$ 的重数为 2、嵌入维数为 3。适当重取解析坐标后，$X$ 由如下形式的方程定义：
\[ z^2+y^3+yq_4(x)+q_6(x)=0\qquad(\mult_0q_i\ge i). \]
$w$ 是由 $w(x,y,z)=(1,2,3)$ 确定的权。
\end{enumerate}
设 $g:\bar Y:=B_0^wX\to X$ 是以权 $w$ 作的加权爆破。此时 $\bar Y$ 只含 Du Val 奇点。$\bar Y$ 被 $Y$ 支配（$h:Y\to\bar Y$），且满足 $K_{\bar Y}\sim\cO_{\bar Y}(1)\sim-h_*Z$。若 $k=2$，则 $B_0^wX$ 与标准爆破 $B_0X$ 的正规化同构。
\end{theorem*}

\begin{remark}{4.58}
上面由权定义的分次依赖于坐标的选取，但理想 $m^w(n)$ 与坐标的选取无关（参见断言 5.37）。
\end{remark}

\begin{proof}
证明分若干阶段进行。

\textbf{阶段 1}\quad $f^*$ 给出自然的同构 $H^0(X,\cO_X)\cong H^0(Y,\cO_Y)$。$Y$ 上的线丛 $L$ 由 $L:=\cO_Y(-Z)\cong\omega_Y$ 确定。对每个 $n$，记
\[ H^0(Y,L^n)=H^0(Y,\cO_Y(-nZ))=:I_n\subset H^0(X,\cO_X), \]
% [原文 p.180]
这些 $I_n$ 构成理想。
由命题 4.45(1) 与命题 4.54(1)，$L^2$ 是 $f$-自由的。由此得到的一般事实“存在适当的 $a\in\Z_{>0}$ 使 $I_{n+a}=I_n\cdot I_a$（$n\gg0$）”在后面的证明中要用到。这里给出更直接的论证：证明 $I_{n+2}=I_n\cdot I_2$（$n>2$）。取 $L^2$ 的两个一般截面 $s,t$，得到正合列
\[ 0\to L^{n-2}\xrightarrow{{}^{t}(t,-s)}L^n\oplus L^n\xrightarrow{(s,t)}L^{n+2}\to0. \]
若 $n>2$，由命题 4.45(2) 与命题 4.51(1) 有 $H^1(L^{n-2})=0$。因此 $I_{n+2}=I_n\cdot(s,t)$，断言得证。

注意 $H^1(Y,\cO_Y(-(n+1)Z))=0$（$n\ge0$），由正合列
\[ 0\to H^0(Y,\cO_Y(-(n+1)Z))\to H^0(Y,\cO_Y(-nZ))\to H^0(Z,L^n|_Z)\to0 \]
得到相应分次环的同构
\[ \bigoplus_{n\ge0}I_n/I_{n+1}\cong\bigoplus_{n\ge0}H^0(Z,L^n|_Z). \]
后者的环结构已在命题 4.54 中给出。下面就利用它来获取关于 $H^0(X,\cO_X)$ 的信息。

由 $H^0(Z,\cO_Z)=\C$ 可知 $f^*m_{0,X}=H^0(Y,\cO_Y(-Z))$。这意味着 $m_{0,X}^n\subset I_n$，但等式一般并不成立。

\textbf{阶段 2}\quad 假设 $k\ge3$。此时由命题 4.54(2)，$\bigoplus I_n/I_{n+1}$ 由 $I_1/I_2$ 生成。因此由阶段 1 得
\[ I_n=I_1^n+I_{n+1}=\cdots=I_1^n+I_{n+4}\subseteq I_1^n+I_n\cdot I_1. \]
由中山引理得 $I_1^n=I_n$，故 $m_{0,X}^n=I_n$ 对所有 $n\ge0$ 成立。

这表明 $m_{0,X}$ 的生成元个数等于 $h^0(Z,L|_Z)=k$。$X$ 的重数为 $k$ 也可由下式看出：
\[ \dim_{\C}(m_{0,X}^n/m_{0,X}^{n+1})=h^0(Z,L^n|_Z)=kn. \]

\textbf{阶段 3}\quad 若 $k\le2$，则 $\bigoplus_{n\ge0}I_n/I_{n+1}$ 存在 3 个生成元 $x,y,z$。仍用同样字母表示它们在 $H^0(X,\cO_X)$ 中的提升。与阶段 2 同样，
% [原文 p.181]
可知 $x,y,z$ 生成 $m_{0,X}\cong I_1$，故 $X$ 的嵌入维数至多为 3。忽略高次项后方程组也提升为 $X$ 的方程组，故用 4.24 的手法即可得到应证的正规形。$I_n=m^w(n)$ 也同样由阶段 2 的方法证明。
% [Tutor Note: 此处 $m_{0,X}\cong I_1$ 中的“$\cong$”疑为“$=$”之误：阶段 2 已证 $k\ge 3$ 时 $m^n=I_n$，阶段 3 只是把该等式移到 $k\le 2$ 的情形，两者应为理想相等而非同构。]

\textbf{阶段 4}\quad 记 $B^*X:=\Proj_X\bigoplus_{n\ge0}f_*\cO_Y(nK_Y)$。由于 $\cO_Y(nK_Y)$（$n\ge3$）由全局截面生成，态射 $p:Y\to B^*X$ 由 $\cO_Y(nK_Y)$（$n\gg1$）的全局截面得到。由 $K_{B^*X}=p_*K_Y$ 与 $K_Y\equiv p^*K_{B^*X}$，可知 $B^*X$ 只含 Du Val 奇点。由 $\Proj$ 的构造还可知 $K_{B^*X}=\cO_{B^*X}(1)$。

若 $k\ge3$，则 $f_*\cO_Y(nK_Y)=m_{0,X}^n$，从而 $B^*X=B_0X$。至此 (1) 的证明完成。

\textbf{阶段 5}\quad 对 $k=1,2$ 的情形已证明 $m^w(n)=f_*\cO_Y(nK_Y)$。因此 $B^*X=\Proj_X\bigoplus_{n\ge0}m^w(n)$。$k=2$ 时还容易看出 $B^*X$ 是标准爆破的正规化。
\end{proof}
% [Tutor Note: 定理 4.57 是定理 4.48 之后的“逆向”构造：椭圆奇点经一次加权爆破即可化为只含 Du Val 奇点者，而 $k=1,2,3$ 恰对应嵌入维数 3、3、k 三种标准形。这也说明椭圆奇点比 Du Val 奇点只多一个离散不变量 $k$。]

\section{超曲面奇点的形变}

\takeaway{本节构造孤立超曲面奇点的半普遍形变空间：把方程写成 $f+\sum t_ig_i$ 的形式，用 Artin 逼近与 Tougeron 引理逐步消去高次项，并讨论解析、形式、代数三种设定下的差别。}

本节目的是构造孤立超曲面奇点的半普遍（miniversal）形变空间。这一结果在本书中只用于 Du Val 奇点的情形，但证明在一般设定下同样成立。有关它的详细讨论可参阅 [Art76, Loo84, AGZV85]。

\begin{definition}{4.59}
设 $0\in X_0$ 是复解析空间的芽。$0\in X_0$ 的形变（deformation）是指复解析空间芽的平坦态射 $f:(0\in X)\to(0\in S)$ 与同构 $f^{-1}(0)\cong X_0$ 组成的对。常省略同构，直接把 $f$ 称为形变。

$0\in X_0$ 的形变 $u:(0\in\mathbf X)\to(0\in\mathbf U)$ 在满足下述条件时称为 $(0\in X_0)$ 的准普遍形变（versal deformation）。设 $f:(0\in X)\to(0\in S)$ 是 $0\in X_0$ 的任意形变。此时存在态射 $u(f):(0\in S)\to(0\in\mathbf U)$，使得 $f:(0\in X)\to(0\in S)$
% [原文 p.182]
同构于 $u$ 沿 $u(f)$ 的拉回
\[
\begin{tikzcd}
X \arrow[r,"\cong"] \arrow[d,"f"'] & S\times_{\mathbf U}\mathbf X \arrow[r] \arrow[d] & \mathbf X \arrow[d,"u"] \\
S \arrow[r,equals] & S \arrow[r,"u(f)"'] & \mathbf U
\end{tikzcd}
\]
且该同构 $X\cong S\times_{\mathbf U}\mathbf X$ 所诱导的同构 $f^{-1}(0)\cong u^{-1}(0)$ 与给定的同一 $f^{-1}(0)\cong X_0$、$u^{-1}(0)\cong X_0$ 一起构成下面的交换图。
\[
\begin{tikzcd}
f^{-1}(0) \arrow[r,"\cong"] \arrow[d] & u^{-1}(0) \arrow[d] \\
X_0 \arrow[r,equals] & X_0
\end{tikzcd}
\]
若进一步，切映射 $du(f):T_0S\to T_0\mathbf U$ 由 $f$ 唯一确定，则称 $u:(0\in\mathbf X)\to(0\in\mathbf U)$ 为半普遍形变（miniversal deformation）\footnote{有的文献把附加了切映射条件的概念直接称为 versal deformation。例如《岩波数学辞典》就把这一受限概念称为 versal deformation，并译作“半普遍形变”。此处与日文译名保持一致。}。此时特别地，不存在解析曲线 $0\in D\subset\mathbf U$ 使 $\mathbf X|_D$ 同构于直积 $X_0\times D$。
\end{definition}

\begin{theorem*}{4.60}
（[Gra72]）对任意孤立奇点，半普遍形变空间存在，且在局部解析同构的意义下唯一确定。
\end{theorem*}

一般情形下，确定半普遍形变空间非常困难。然而孤立超曲面奇点的半普遍形变空间可以具体写出。

\begin{theorem*}{4.61}
设 $0\in X_0=(f(x_1,\dots,x_m)=0)\subset\C^m$ 是原点处孤立超曲面奇点的芽。选取收敛幂级数 $g_1,\dots,g_n$（$n=\tau(X_0)$），使其成为
\[ \cO_{0,\C^m}/(f,\partial f/\partial x_1,\dots,\partial f/\partial x_m) \]
的基。于是
% [原文 p.183]
\[
\begin{tikzcd}
\mathbf X:=\Bigl\{f(x_1,\dots,x_m)+\sum_i t_ig_i(x_1,\dots,x_m)=0\Bigr\} \arrow[d] \arrow[r,hook,"\subset"] & \C^m\times\C^n \arrow[d] \\
\mathbf U:=\C^n \arrow[r,equals] & \C^n
\end{tikzcd}
\]
是 $X_0$ 的半普遍形变。
\end{theorem*}
% [Tutor Note: 这里的 $n=\tau(X_0)$ 是 Tjurina 数，即 $\cO_{0,\C^m}/(f,\partial f)$ 的维数；它与 Milnor 数 $\mu=\dim\cO/(\partial f)$ 不同，对拟齐次奇点二者相等，一般情形 $\tau\le\mu$。]

作为证明的第一步，先把断言改写成方程的语言。也就是说，定理 4.61 可由下述结论得到。

\begin{theorem*}{4.62}
记号同上。设 $(0\in T)$ 是复解析空间的芽，$F(x,t)\in\cO_{0,\C^m\times T}$ 是满足 $F(x,0)=f(x)$ 的收敛幂级数。此时存在 $a_j(t)\in(t)\cO_{0,T}$（$j=1,\dots,n$）与 $b_i(x,t),d(x,t)\in(t)\cO_{0,\C^m\times T}$（$i=1,\dots,m$），满足下式：
\[ (1+d(x,t))\cdot F(x,t)=f(x-b)+\sum_j a_j(t)g_j(x-b). \]
其中 $x={}^t(x_1,\dots,x_m)$，$b={}^t(b_1,\dots,b_m)$。此外，$a_j(t)$ 模 $(t)^2\cO_{0,T}$ 是唯一的。
\end{theorem*}

证明的主要部分由两个引理组成。第一个引理在“忽略高次项”的意义下证明定理 4.62，第二个引理则断言“充分高次的项”可以去掉。

\begin{lemmap}{4.63}
记号同上。设 $J\subset\cO_{0,\C^m\times T}$ 是使 $\cO_{0,\C^m\times T}/J$ 在 $\cO_{0,T}$ 上有限的理想。此时存在收敛幂级数 $a_j(t)\in(t)\cO_{0,T}$ 与 $b_i(x,t),d(x,t)\in(t)\cO_{0,\C^m\times T}$，满足
\[ (1+d(x,t))\cdot F(x,t)\equiv f(x-b)+\sum_j a_j(t)g_j(x-b)\mod J\cO_{0,\C^m\times T}. \]
\end{lemmap}

\begin{proof}
取 $v_k\in\cO_{0,\C^m\times T}$（$k=1,\dots,M$）为 $\cO_{0,T}$-模 $\cO_{0,\C^m\times T}/J$ 的生成元的代表，记 $v={}^t(v_1,\dots,v_M)$。引入以独立变量 $\eta_k$、$\xi_{i,k}$ 为元素的向量 $\eta=(\eta_1,\dots,\eta_M)$ 与 $m\times M$ 矩阵 $\xi=(\xi_{i,k})$。

为简化记号，记 $f^+=f(x-\xi\cdot v)$、$f_i^+=\partial f/\partial x_i(x-\xi\cdot v)$、$g_j^+=g_j(x-\xi\cdot v)$；再把 $\xi,\eta$ 作为变量添入，记 $\cO=\cO_{0,\C^{m+mM+M}\times T}$、$\cO^+=\cO_{0,\C^{mM+M}\times T}$。
% [原文 p.184]
首先证明自然态射
\[ f^+\cdot(\cO/J\cO)\oplus\bigoplus_i f_i^+\cdot(\cO/J\cO)\oplus\bigoplus_j g_j^+\cdot\cO^+\to\cO/J\cO \]
是满射。由关于 $J$ 的假设，$\cO/J\cO$ 在 $\cO^+$ 上有限。因此由中山引理，只需在点 $0\in\C^{mM+M}\times T$ 上（即 $t=\eta=\xi=0$ 时）验证它是满射，而这一情形是显然的。

用 $(t,\eta,\xi)$（或 $(t)$）表示 $\cO^+$（或 $\cO_{0,T}$）的极大理想。此时
\[ (1+\eta\cdot v)\cdot F(x,t)-(1+\eta\cdot v)f^+-(f_1^+,\dots,f_m^+)\cdot\xi\cdot v\in(t)\cO+(t,\eta,\xi)^2\cO \]
成立。这可由假设 $F(x,0)=f(x)$ 与下面的 Taylor 公式得到：
\[ f^+=f(x-\xi\cdot v)\equiv f(x)-(f_1,\dots,f_m)\cdot\xi\cdot v\mod(\xi_{i,k})^2. \]

由上述两个断言，存在以 $\xi,\eta,t$ 为变量的收敛幂级数 $A_j$、$U_k$、$W_{i,k}$，使得模 $(t)\cO^++(t,\eta,\xi)^2\cO^+$ 有 $A_j\equiv0$、$U_k\equiv\eta_k$、$W_{i,k}\equiv\xi_{i,k}$，且满足
\[ (1+\eta\cdot v)F(x,t)\equiv\Bigl(1+\sum_kU_kv_k\Bigr)f^++\sum_{i,k}f_i^+W_{i,k}\cdot v_k+\sum_jA_jg_j^+\mod J\cO_{0,\C^{m+mM+M}\times T}. \]
要看出这一点，只需对展开式 $(1+\eta\cdot v)F(x,t)-\cdots=\sum_it_ip_i+\sum_{j,k}\eta_j\eta_kq_{j,k}+\cdots$ 中出现的 $p_i,q_{j,k},\cdots\in\cO$ 代入模 $J$ 的表示。

幂级数 $U_k$ 与 $W_{i,k}$ 关于变量 $\xi_{i,k},\eta_k$ 含有独立的线性项，常数项属于 $(t)\cO_{0,T}$。因此由隐函数定理，方程组 $U_k=0$、$W_{i,k}=0$（$\forall i,k$）有唯一解 $\eta_k(t),\xi_{i,k}(t)\in(t)\cO_{0,T}$。令 $a_j(t)=A_j(t,\eta(t),\xi(t))$、$b=\xi(t)\cdot v$、$d=\eta(t)\cdot v$ 即可。
\end{proof}
% [Tutor Note: 引理 4.63 的实质是把“模高次项求形变”问题线性化：由 $J$ 的有限性保证 $\cO/J\cO$ 在 $\cO^+$ 上有限，于是中山引理把满射性归结到原点一个点的线性代数计算。]

第二个结果本质上是以 Tougeron 引理之名在 [Art69a] 中介绍的。

\begin{lemmap}{4.64}
记号同上。设 $I=(F,\partial F/\partial x_1,\dots,\partial F/\partial x_m)$，$I_1\subset(x,t)\cO_{0,\C^m\times T}$ 是任意理想，$G(x,t)\in\cO_{0,\C^m\times T}$ 是满足
\[ F(x,t)\equiv G(x,t)\mod I^2\cdot I_1 \]
的任意元。此时存在收敛幂级数 $b_i(x,t),d(x,t)\in I_1$，使得
% [原文 p.185]
\[ (1+d(x,t))\cdot F(x,t)=G(x-b,t). \]
\end{lemmap}

\begin{proof}
为简便，用 $F_i$ 表示 $\partial F/\partial x_i$。选取 $c\in I\cdot I_1$ 与 $c_{i,i'}\in I_1$，使 $F(x,t)-G(x,t)=\sum_{i,i'}c_{i,i'}F_iF_{i'}+c\cdot F$ 成立。把 $c_{i,i'}$ 换成 $c_{i,i'}/(1-c)$，即可设 $F(x,t)-G(x,t)/(1-c)=\sum_{i,i'}c_{i,i'}F_iF_{i'}$。记 $\nabla F:={}^t(F_1,\dots,F_m)$，并设 $\xi:=(\xi_{i,i'})$ 是以变量为元素的 $m\times m$ 矩阵。由 Taylor 公式有
\[ F(x+\xi\cdot\nabla F,t)\equiv F(x,t)+\sum_{i,i'}\xi_{i,i'}F_iF_{i'}\mod(\xi_{i,i'})^2(F_j)^2. \]
因此存在关于 $x,\xi$ 的收敛幂级数 $W_{i,i'}$，使得
\[ W_{i,i'}\equiv c_{i,i'}+\xi_{i,i'}\mod(\xi_{i,i'})^2\cO_{0,\C^{m+m^2}\times T}, \]
\[ F(x+\xi\cdot\nabla F,t)=G/(1-c)+\sum_{i,i'}W_{i,i'}F_iF_{i'}. \]
函数 $W_{i,i'}$ 关于变量 $\xi_{i,i'}$ 含有独立的线性项，常数项属于 $I_1$。由隐函数定理，方程组 $W_{i,i'}=0$（$\forall i,i'$）有唯一解 $\xi_{i,i'}(x,t)\in(x,t)\cO_{0,\C^m\times T}$。再由 $W_{i,i'}$ 的形状得到理想的包含关系 $(\xi_{i,i'})\subset(\xi_{i,i'})^2+I_1$，由中山引理即得 $\xi_{i,i'}\in I_1$。于是令 $b={}^t(b_1,\dots,b_m)=\xi(x,t)\cdot\nabla F$，则 $b_i\in I_1$，且
\[ (1-c(x,t))\cdot F(x+b,t)=G(x,t). \]
把 $x$ 换成 $x-b$ 得 $(1-c(x-b,t))\cdot F(x,t)=G(x-b,t)$。由 $b_i\in I_1$ 得 $c(x-b,t)\in I_1$。
\end{proof}

\begin{proof}[4.65（定理 4.62 的证明）]
记 $I=(F,\partial F/\partial x_1,\dots,\partial F/\partial x_m)$。由于
\[ \cO_{0,\C^m\times T}/I+(t)\cong\C\{x\}/(f,\partial f/\partial x_1,\dots,\partial f/\partial x_m) \]
是 Artin 环，故对任意 $\ell>0$，$\cO_{0,\C^m\times T}/I^\ell+(t)$ 也是 Artin 环。因此 $\cO_{0,\C^m\times T}/I^\ell$ 在 $\cO_{0,T}$ 上有限。取 $\ell\gg0$，由 Artin-Rees 引理 [Mat86, 8.5] 有 $(t)\cap I^\ell\subset(t)I^2$。首先把引理 4.63 用于 $J=I^\ell$，并令 $G(x,t)=f(x-b)+\sum_ja_j(t)g_j(x-b)$，则 $(1+d)F-G\in(t)\cap I^\ell\subset(t)I^2$。因此令 $I_1=(t)$，对 $(1+d)F$ 与 $G$ 应用引理 4.64 即可。（注意把 $F$ 换成 $(1+d)F$ 时 $I$ 不变。）
% [原文 p.186]
其次考察唯一性。模 $((t)^2+(t)I)\cO_{0,\C^m\times T}$ 来看，注意 $t_i\cdot F\equiv0$ 以及 $f(x-b)\equiv f(x)$ 等，得到
\[ F(x,t)\equiv f(x)+\sum_ja_j(t)g_j(x)\mod((t)^2+(t)I)\cO_{0,\C^m\times T}. \]
对 $t_i$ 求偏导数，注意 $(t)+I=(t)+(f,f_1,\dots,f_m)$，得到
\[ \frac{\partial F}{\partial t_i}(x,0)\equiv\sum_j\frac{\partial a_j}{\partial t_i}(0)g_j(x)\mod(f,f_1,\dots,f_m). \]
$a_j$ 模 $(t)^2\cO_{0,T}$ 的唯一性由此显然。
\end{proof}

\par\medskip\noindent\textbf{4.66（代数情形下的结果）}\footnote{以下内容与“尽量不使用开簇、解析空间等来构建射影簇的极小模型纲领理论”这一工作相关。若无兴趣，此后本节可以跳过。}\quad 本节一直在解析的设定下讨论，但其他情形下论证同样有效。首先，形式幂级数的情形不需要对论证作任何改动。

更重要的是理解在代数设定下应当如何处理。若 $f(x)$、$F(x,t)$ 等对象是代数簇芽上的正则有理函数（regular function）\footnote{为避免与分析函数混淆，以下本节把代数簇情形下的 regular function 称为正则有理函数。}，则引理 4.64 中的 $W_{i,i'}$ 也可以是（在芽的充分小的适当代表上的）正则有理函数。唯一的困难在于，隐函数定理对正则有理函数并不成立。也就是说，多数情形下解 $\xi_{i,i'}(x,t)$ 不是芽的任何代表上的正则有理函数。（用传统术语来说，$\xi_{i,i'}(x,t)$ 是 $\C^m\times T$ 上的代数函数。参见 [PS97, Sie69]。）

从代数空间（algebraic space）（参见 [Art69b, Art70]）的观点看，要点在于：在方程组 $W_{i,i'}=0$ 所定义的代数簇 $T'$ 上，$\xi_{i,i'}$ 是正则有理函数，
% [原文 p.187]
且到 $\C^m\times T$ 的投影在原点处是 étale 的（定义 4.67）。

同样地，若 $f$ 是（由正则有理函数组成的）局部环 $\cO_{0,\C^m}$ 的元，则为使引理 4.63 成立，必须作 étale 的基变换 $(0\in T')\to(0\in T)$。引理 4.64 所断言的是：两个超曲面拉回到 $\C^m\times T'$ 上的 étale 空间后，会成为代数同构的。

下面陈述定理 4.61 的代数类似（命题 4.69），以说明上述论证实际上应当如何修正。除由复解析条件导出的推论 4.68 之外，所有结果在任意特征的代数闭域上都成立。关于代数形变理论的更详细解说可参阅 [Art76]。\hfill$\square$

\begin{definition}{4.67}
设 $f:(\xi\in X)\to(\eta\in Y)$ 是 $\C$ 上代数概型芽之间的态射。

若 $(\xi\in X)$ 具有以 $x={}^t(x_1,\dots,x_n)$ 为变量的如下表示
\[ \cO_{\xi,X}=(\cO_{\eta,Y}[x])_{\xi}/(g_1,\dots,g_n), \]
其中 $\xi=(0,\eta)\in\C^n\times Y$、各 $g_i$ 满足 $g_i(\xi)=0$，且 Jacobi 矩阵 $\partial g/\partial x(\xi):=(\partial g_i/\partial x_j(\xi))$ 可逆，则称 $f$ 在 $\xi$ 处是 étale（平展）的。

若 $f$ 是 étale 的，则隐函数定理意味着 $\cO_{\xi,X}^{an}\cong\cO_{\eta,Y}^{an}$。特别地，$(\xi\in X)$ 的嵌入维数 $\dim_{\C}m_{\xi,X}/m_{\xi,X}^2$ 与 $\xi$ 处的维数 $\dim_{\xi}X$，分别等于 $(\eta\in Y)$ 的对应维数。

给定两个芽 $(\xi\in X)$ 与 $(\eta\in Y)$，若存在芽 $(\zeta\in Z)$ 与两个 étale 态射 $(\zeta\in Z)\to(\xi\in X)$、$(\zeta\in Z)\to(\eta\in Y)$，则称这些芽是 étale 等价的，记作 $(\xi\in X)\cong_{et}(\eta\in Y)$。
\end{definition}

\begin{corollary}{4.68}
（引理 4.64 及其代数类似的推论）设给定两个在原点处有孤立奇点的多项式超曲面（含于 $(0\in\C^n)$）。它们 étale 等价，当且仅当它们复解析同构。
\end{corollary}

\begin{propositionp}{4.69}
符号同定理 4.61。设 $f,g_1,\dots,g_n$ 是正则有理函数。设 $f:(0\in X)\to(0\in T)$ 是带点代数概型的平坦态射，并给定同一 $(0\in f^{-1}(0))\cong_{et}(0\in X_0)$（参见推论 4.68）。此时适当选取概型的 étale 态射 $(0'\in T')\to(0\in T)$，可使诱导的形变 $f':(0'\in X'=X\times_TT')\to(0'\in T')$
% [原文 p.188]
满足：（定义 4.59 中的）$u(f'):T'\to\mathbf U$ 是概型的态射。
\end{propositionp}

\begin{proof}
下面说明定理 4.62 的证明应如何修正。

可设 $0\in X$ 是奇点。设 $(\zeta\in Z)$ 同时具有到 $(0\in f^{-1}(0))$ 与 $(0\in X_0)$ 的 étale 态射。

由推论 4.71 有 $(\zeta\in Z)\subset(0\in\C^m)$。把态射 $(\zeta\in Z)\to(0\in f^{-1}(0))$ 扩张为态射 $(0\in\C^m)\to(0\in\C^m)$。该态射必然是 étale 的。（对 $X_0$ 也作同样处理。）于是可以把一切拉回到 $(0\in\C^m\times T)\supset(\zeta\in Z)\times\{t\}$ 上，故可设 $(0\in f^{-1}(0))=(0\in X_0)$。

设 $F(x,t)=0$ 是 $0$ 附近 $X\subset\C^m\times T$ 的方程。令
\[ I=(F,\partial F/\partial x_1,\dots,\partial F/\partial x_m). \]
由于 $\Spec\cO_{0,\C^m\times T}/(I+(t))$ 有限，概型 $\Spec\cO_{0,\C^m\times T}/I$ 在 $0\in T$ 上拟有限。由引理 4.72，存在 $0\in T$ 上的 étale 带点概型 $0'\in T'$，使诱导态射 $\Spec\cO_{0',\C^m\times T'}/I\cO_{0',\C^m\times T'}\to T'$ 成为有限态射。作这样的基变换，并把 $T'$ 改写为 $T$，即可沿用与定理 4.62 相同的记号。

取多项式 $v(x,t)$、$u(x)$，使 $F\cdot v$、$f\cdot u$、$g_1\cdot u,\dots,g_n\cdot u$ 都是多项式，且 $v(0,0)u(0)\ne0$、$F(x,0)\cdot v(x,0)=f\cdot u$ 成立。把 $f$ 换成 $u\cdot f$、$g_j$ 换成 $u\cdot g_j$，理想 $(f,\partial f/\partial x_1,\dots,\partial f/\partial x_m)$ 不变，$g_j$ 仍是商向量空间的基。把 $F$ 换成 $F\cdot v$，理想 $I$ 也不变。这样便可设 $F,f,g_1,\dots,g_n$ 都是多项式且 $F(x,0)=f(x)$。

这里把引理 4.63 的论证用于 $J=I^\ell$（$\ell\gg0$），并沿用其记号。注意 $U_k$、$W_{i,k}$、$A_j$ 都是多项式。由 $U_k=W_{i,k}=0$（$\forall i,k$）定义的 $\C^{mM+M}\times T$ 的闭子概型 $(0'\in T')$ 到 $(0\in T)$ 的投影是 étale 的。$\eta_k$、$\xi_{i,k}$ 在 $T'$ 上的限制满足 $U_k=W_{i,k}=0$。因此存在正则有理函数 $a'_j\in(t)\cO_{0',T'}$、$b'_j,d'\in(t)\cO_{0',\C^m\times T'}$，满足
% [原文 p.189]
\[ (1+d')\cdot F(x,t')\equiv f(x-b')+\sum_ja'_jg_j(x-b')\mod J\cO_{0,\C^m\times T'}. \]
由引理 4.64，$(0'\in F(x,t')=0)\subset\C^m\times T'$ 在 $(0'\in T')$ 上与 $(0\in f(x-b')+\sum_ja'_jg_j(x-b')=0)\subset\C^m\times T'$ 复解析同构。由此可知 $(a'_1,\dots,a'_n):T'\to\mathbf U$ 是概型的态射。
\end{proof}

下面说明上面用到的辅助结果。

\begin{propositionp}{4.70}
设 $f:(x\in X)\to(t\in T)$ 是 $\C$ 上代数概型芽的态射，$d=\dim_xf^{-1}(t)$。此时 $f$ 可分解为开嵌入（open embedding）$g:(x\in X)\to Y$ 与满足 $d=\dim h^{-1}(t)$ 的射影态射 $h:Y\to T$。
\end{propositionp}

\begin{proof}
对适当的 $N$，存在从芽 $(x\in X)$ 到闭子概型 $X_1\subset\PP^N\times T$ 的开嵌入。用 $d+1$ 个一般超平面 $H_0,\dots,H_d$ 构成一般线性投影 $\pi:\PP^N\times T\dashrightarrow\PP^d\times T$，则它在 $x$ 处是态射。爆破闭子概型 $Z=\bigcap_iH_i$，便得到芽 $(x\in X)$ 到 $X_2=B_ZX_1$ 的开嵌入与射影态射 $\phi:X_2\to\PP^d\times T$。设 $g:X_2\to Y$ 是 $\phi$ 的 Stein 分解 [Har77, III.11.5]。由于 $Y\to\PP^d\times T$ 是有限态射，诱导态射 $Y\to\PP^d\times T\to T$ 是射影的，且纤维维数为 $d$。在 $g(x)$ 的邻域上，$g$ 是有限态射（特别地是仿射的），且 $g_*\cO_{X_2}=\cO_Y$ 成立。于是 $g$ 在 $x$ 的邻域上是同构。
\end{proof}

\begin{corollaryp}{4.71}
（极小嵌入（minimal embedding））符号同上。设 $m$ 是 $f^{-1}(t)$ 在 $x$ 处的嵌入维数。此时存在代数概型芽的嵌入 $(x\in X)\subset((0,t)\in\C^n\times T)$，其中 $n=\max\{m,d+1\}$。
\end{corollaryp}

\begin{proof}
由命题 4.70，选取嵌入 $X\subset\PP^N\times T$（$N>0$），使其闭包 $\bar X$ 在 $t$ 上具有纤维维数 $d$。此时一般线性投影 $g:\bar X\to\PP^d\times T$ 是有限态射。设 $x_i$ 是 $g^{-1}(g(x))$ 中异于 $x$ 的点。这里取一般线性投影 $h:\bar X\to\PP^n\times T$，使其经过 $g$ 分解。由 $n>d$ 得 $h(x_i)\ne h(x)$（$\forall i$），又由 $n\ge m$ 得 $x$ 处的相对切映射 $dh(x)$ 是单射。因此 $h$ 在 $x$ 的邻域上是嵌入。
\end{proof}

% [原文 p.190]
\begin{lemmap}{4.72}
设 $(x\in X)\to(y\in Y)$ 是 $\C$ 上代数概型芽的拟有限态射。此时可适当选取 étale 态射 $(y'\in Y')\to(y\in Y)$，使得纤维积 $X\times_YY'$ 中存在既开又闭的集合 $X'\ni(x,y')$，并且 $X'$ 在 $Y'$ 上有限、$(x,y')$ 是 $y'$ 上的唯一一点。
\end{lemmap}

\begin{proof}
应用命题 4.70 的 $d=0$ 情形，可把态射 $(x\in X)\to(y\in Y)$ 分解为开嵌入 $(x\in X)\subset Z$ 与有限态射 $f:Z\to Y$。

设 $v_1,\dots,v_n\in\cO_Z$ 是 $\cO_Y$-模 $\cO_Z$ 的生成元。记 $\Sigma=f^{-1}(y)$。此时 $\cO_Z/m_{y,Y}\cO_Z=\bigoplus_{z\in\Sigma}\cO_{z,Z}/m_{y,Y}\cO_{z,Z}$，故可适当选取 $u\in\cO_Z$，使 $u\in m_{y,Y}\cO_{x,Z}$ 与 $u-1\in m_{y,Y}\cO_{z,Z}$（$\forall z\in\Sigma\setminus\{x\}$）同时成立。特别地 $u^2-u\in m_{y,Y}\cO_Z$。取 $a_i\in\cO_{y,Y}$ 使 $u=\sum_ia_iv_i$ 成立。以 $\xi=(\xi_1,\dots,\xi_n)$ 为变量，引入多项式 $\phi(\xi):=\sum_i(\xi_i+a_i)v_i\in\cO_Z[\xi]$。由于 $\phi(\xi)=\sum_i\xi_iv_i+u$，有
\[ \phi^2-\phi\equiv\Bigl(\sum_i\xi_iv_i\Bigr)(2u-1)\mod m_{y,Y}\cO_Z+(\xi)^2\cO_Z. \]
由 $(2u-1)^2\equiv1\mod m_{y,Y}\cO_Z$ 知 $2u-1\in\cO_Z$ 是可逆元。因此可选取多项式 $A_i\in\cO_{y,Y}[\xi]$，使 $\phi^2-\phi=(2u-1)\sum_iA_iv_i$ 与 $A_i-\xi_i\in m_{y,Y}\cO_{y,Y}[\xi]+(\xi)^2$ 成立。（$\{A_i\}$ 的取法不唯一，任取一组即可。）

于是 $Y'=\Spec\cO_{y,Y}[\xi]/(A_1,\dots,A_n)\to Y$ 在 $y'=(y,0)$ 处是 étale 的。用 $\xi'_i$ 表示 $\xi_i$ 的像，则在 $\cO_Z\otimes_{\cO_{y,Y}}\cO_{y',Y'}$ 中 $\phi'=\sum_i\xi'_iv_i+u$ 是幂等元，且 $X'=(\phi'=0)$ 在 $y'$ 上只含 $(x,y')$。因此 $X'\subset X\times_YY'$。
\end{proof}

\chapterreview{第 4 章在全书中的角色，是提供极小模型纲领的“二维局部模型”。纲领的真正舞台在三维及以上：在那里，终端奇点、标准奇点与 klt 奇点取代了光滑性，而这些概念的二维原型可以彻底分类。本章把二维情形一次做完——既给出分类，也给出把奇点化为 Du Val 奇点的构造，因此第 5 章可以把它当作现成的局部结论引用（例如定理 5.35 处理三维标准奇点时就用到椭圆奇点）。

在知识脉络上，本章接续第 1 章的奇点消解与差异、第 2 章的伴随与消没定理、第 3 章的对数终端语言：记号 4.44 中 $K_Y\equiv\sum a_iE_i$、$Z=\sum[-a_i]E_i$ 的分解正是这些工具的直接产物。反过来，本章为后文铺路：Du Val 奇点的同时奇点消解（定理 4.28 及其整体化定理 4.43）是 DLT 形变与锥定理局部化的技术范本；超曲面奇点的半普遍形变理论则提供了形变论意义上的通用工具箱。

核心结论可分五组。§4.1 讨论二维对数标准奇点，说明它与商奇点的关系；§4.2 给出 Du Val 奇点的 $A_n$、$D_n$、$E_6$、$E_7$、$E_8$ 分类及例外曲线图；§4.3 证明 Du Val 奇点族的纤维可以同时消解，并在有限基变换后整体实现；§4.4 把椭圆奇点归结为三条标准方程——$k=-(Z\cdot Z)\ge3$ 时是 $\mathbb{A}^k$ 中的超曲面，$k=2$ 时是 $z^2+q_4(x,y)=0$，$k=1$ 时是 $z^2+y^3+yq_4(x)+q_6(x)=0$——并用加权爆破一步化为 Du Val 奇点；§4.5 证明孤立超曲面奇点的半普遍形变空间存在，并可由 $f+\sum t_ig_i$ 显示写出。

有三点需要辨析。第一，除子 $Z$ 有两种来源：§4.4 的记号 4.44 用 $K_Y\equiv\sum a_iE_i$ 定义 $Z=\sum[-a_i]E_i$，而定理 4.57 用 $\omega_Y=f^*\omega_X(-Z)$ 定义；当 $K_X$ 是 Cartier 时二者一致，否则后者只是一个数值 $f$-平凡的例外除子，不能与前者混同。第二，$k=-(Z\cdot Z)$ 越大，嵌入维数越大，这与“自交数为负”的直觉相反，源头在于 $Z$ 的系数为整部而非分数部。第三，证明紧跟其后的陈述不带尾符、证明在别处者带 □、书中不证者带 ■，这是本书特有的约定；读时不要以符号判断命题的重要性。

从历史看，Du Val 奇点的 $A$–$D$–$E$ 分类源于 Du Val 1934 年的工作，Brieskorn 用它给出了同时消解的经典构造，Artin 则把逼近与形变论引入这一领域；半普遍形变空间的存在性是 Grauert 的定理，而定理 4.61 的具体描述与 Tjurina 数 $\tau$ 联系在一起。椭圆奇点的现代处理属于 Laufer 与 Reid。读法上，§4.4 的三条标准方程最好与 §4.3 的 Du Val 分类对照着看，才能体会“多一个离散不变量 $k$”的含义；§4.5 技术性强且后文用量不大，按脚注提示可以跳读，只记住“超曲面奇点的半普遍形变可由 $f+\sum t_ig_i$ 显式给出”这一结论即可。}
